%!TeX root=../main.tex

\chapter{A escolha do OpenTelemetry}

\label{cap:opentelemetry}

O caminho que falta ao kw tem um componente central, o OpenTelemetry. Este
capítulo explica o que ele é e justifica a escolha, comparando-o com
alternativas e com o que outras ferramentas de linha de comando fizeram, antes
que o Capítulo~\ref{cap:solucao} mostre como ele é usado.

\section{O que é o OpenTelemetry}
\label{sec:otel-o-que-e}

Instrumentar um programa para emitir dados de telemetria costuma amarrar o
código à ferramenta de análise que recebe esses dados, porque o formato e o
protocolo são os dela. O OpenTelemetry existe para desfazer esse laço. É um
arcabouço de observabilidade, com ferramentas para gerar, exportar e coletar
dados de telemetria como \textit{traces}, métricas e \textit{logs}, e se
declara independente de fornecedor e de ferramenta: o mesmo dado pode ir para
ferramentas diferentes \citep{otel_what_is}. O projeto surgiu em 2019 da fusão
de dois projetos anteriores, o OpenTracing e o OpenCensus, e entrou na Cloud
Native Computing Foundation (CNCF) no estágio inicial, o \textit{sandbox}
\citep{cncf_otel_history_2019}. Em maio de 2026 chegou ao nível
\textit{graduated} \citep{cncf_otel_project}, que o processo da fundação
descreve como o ápice da maturidade de um projeto
\citep{cncf_project_lifecycle}.

O resto do capítulo usa quatro conceitos, apresentados na ordem em que o dado
anda:
\begin{itemize}
  \item \emph{Sinais}: os tipos de dado que o projeto define, isto é,
    \textit{traces}, que mostram o que acontece quando uma requisição é feita,
    métricas, que são medições tomadas durante a execução, e \textit{logs},
    registros de eventos com carimbo de tempo e metadados opcionais
    \citep{otel_signals}, cada um com o próprio estado de estabilidade (o
    protocolo é estável para os três, e o sinal de \textit{profiles} segue em
    desenvolvimento) \citep{otlp_spec_1_11_1}.
  \item \emph{OTLP}: o protocolo que leva qualquer sinal de quem o emite a
    quem o recebe, com dois transportes, gRPC e HTTP; no HTTP, o corpo da
    requisição pode ser protobuf binário ou JSON \citep{otlp_spec_1_11_1}.
  \item \emph{Collector}: um programa pronto que recebe dados de telemetria,
    os processa e os envia a um ou mais sistemas de armazenamento ou análise, organizado em
    \textit{pipelines} de \textit{receivers}, que recebem o dado,
    \textit{processors}, opcionais, que o transformam, e \textit{exporters},
    que o mandam para fora do Collector \citep{otel_collector_architecture}.
  \item \emph{SDKs}: bibliotecas que um programa usa para produzir e exportar
    sinais, uma para cada linguagem atendida, como a de Python
    \citep{otel_python_docs}.
\end{itemize}

Protocolo e SDK são componentes separáveis. A especificação do OTLP define o que
trafega na rede: no transporte HTTP, o cliente envia os dados numa requisição
POST \citep{otlp_spec_1_11_1}. Qualquer programa que monte a mensagem e a
envie fala OTLP, com SDK ou sem ele. Assim, a escolha de quem recebe os dados,
o Collector, fica separada da escolha de como enviá-los, e as seções seguintes
tratam de cada uma à parte.

\section{Critérios de escolha}
\label{sec:criterios}

O motivo central da escolha é o Collector. Ele é um receptor pronto, que já
servia ao caso e se monta por configuração, a partir de \textit{receivers},
\textit{processors} e \textit{exporters} que já existem
\citep{otel_collector_architecture}; o receptor de OTLP e o exportador por
HTTP são componentes estáveis para \textit{logs} na versão usada
\citep{otel_collector_v0_121_0}. Sem ele, seria preciso ler a especificação e
escrever o servidor e o tratamento do protocolo. O segundo motivo é o uso
amplo. Plataformas de análise de vários fornecedores recebem OTLP, como o
Sentry, o PostHog, o Datadog e o New Relic \citep{sentry_otlp, posthog_logs,
datadog_otlp_agent, newrelic_otlp}, de modo que o dado coletado não fica preso
à plataforma escolhida hoje. A maturidade é um ponto distinto, atestado pelo
nível \textit{graduated} da CNCF \citep{cncf_otel_project}.

Do OpenTelemetry, o trabalho usa só o OTLP sobre HTTP e um Collector
configurado, sem SDK. Quem envia precisa apenas de uma requisição HTTP com
corpo JSON, que um programa feito só com a biblioteca padrão do Python monta.
Do lado do servidor, o Collector assume o receptor voltado aos clientes e o
limite de memória, por configuração, sem código de servidor. Sobra um backend
próprio, pequeno, que entende uma parte do OTLP: recebe o corpo que o Collector
reenvia, separa os registros e põe à parte os malformados, descarta os
repetidos, grava num banco SQLite e confirma. O Capítulo~\ref{cap:solucao}
mostra como o Collector e o backend foram montados.

Os critérios da Tabela~\ref{tab:criterios}, de C1 a C9, vêm dos requisitos dos
Capítulos~\ref{cap:kworkflow} e~\ref{cap:motivacao} e de requisitos de
engenharia, e foram fixados antes de olhar as alternativas. O
OpenTelemetry, porém, não entrou na comparação como um candidato entre
outros. Ele era a hipótese de partida, por causa do reuso do Collector, e a
comparação serviu para conferir se a hipótese se sustentava contra os
critérios.

\begin{table}[!ht]
  \centering
  \caption{Critérios da escolha: requisitos (C1 a C8) e preferências (C9).}
  \label{tab:criterios}
  \footnotesize
  \begin{tabular}{@{}>{\raggedright\arraybackslash}p{4.6cm}
                    >{\raggedright\arraybackslash}p{9.6cm}@{}}
    \toprule
    Critério & O que exige \\
    \midrule
    \multicolumn{2}{@{}l}{\textit{Requisitos}} \\
    \addlinespace[2pt]
    C1 Consentimento e filtro por campo
      & O dado só sai após o opt-in, e a privacidade se aplica
        campo a campo, no cliente, por configuração, sem mudar código.\\
    \addlinespace[1pt]
    C2 Extensibilidade
      & Uma fonte ou um atributo novo vai do cliente até a consulta sem
        mudança no protocolo, no Collector nem no backend.\\
    \addlinespace[1pt]
    C3 Um registro por evento, com o tempo do evento
      & Cada evento vai como um registro, com atributos chave e valor livres e
        o momento em que ocorreu, com granularidade de semana; a agregação fica com o backend.\\
    \addlinespace[1pt]
    C4 Cliente em Bash e em Python só com a biblioteca padrão
      & Sem processo residente nem dependência instalada à parte; cada passo
        traduzível para \texttt{sqlite3}, \texttt{jq} e \texttt{curl}.\\
    \addlinespace[1pt]
    C5 Linux, aberto, auto-hospedável
      & Roda em Linux, tem código aberto e pode ser hospedado por quem recebe
        os dados, sem conta em serviço de terceiros.\\
    \addlinespace[1pt]
    C6 Custo operacional baixo
      & Uma máquina, configuração em vez de código de servidor, sem banco ou
        fila obrigatórios para começar.\\
    \addlinespace[1pt]
    C7 Garantia de entrega
      & Confirmação só depois da gravação no backend; reenvio sem duplicar.\\
    \addlinespace[1pt]
    C8 Maturidade e estabilidade
      & Sinal e protocolo estáveis na especificação, versão fixada e histórico
        de mudanças incompatíveis conhecido.\\
    \midrule
    \multicolumn{2}{@{}l}{\textit{Preferências}} \\
    \addlinespace[2pt]
    C9 Reuso de um receptor pronto e amplamente usado
      & O lado servidor e o tratamento do protocolo vêm prontos e
        configuráveis, em vez de lidos da especificação e escritos pelo
        projeto, e o componente tem uso amplo.\\
    \bottomrule
  \end{tabular}
\end{table}


\section{Alternativas consideradas}
\label{sec:alternativas}

A comparação põe o OpenTelemetry ao lado de três outras formas de montar o
caminho inteiro, do cliente ao backend: não reusar nada e escrever um
receptor próprio, usar um sistema de métricas, ou usar uma plataforma de
análise pronta. Cada uma representa um grupo; outras opções examinadas
ficaram de fora, e a lista não é exaustiva. Aqui se pergunta por que o
OpenTelemetry. Como ele é usado, com que sinal, que cliente e que
codificação, é escolha do Capítulo~\ref{cap:solucao}. Os critérios estão na
Tabela~\ref{tab:criterios}, e a parte seguinte explica as avaliações de cada
alternativa, resumidas na Tabela~\ref{tab:alternativas}, ao fim da seção. As
ressalvas do próprio OpenTelemetry ficam para a Seção~\ref{sec:limites-otel}.


A primeira alternativa é um \textbf{servidor HTTP próprio}, que recebesse JSON
no formato que o kw quisesse. É o caminho mais curto no começo e, no papel,
atende quase todos os critérios. Tem também vantagens sobre a escolha: a
superfície é menor, sem Collector e sem as regras do JSON do OTLP, e o
contrato de erro é do projeto. Seu custo está no C9: seria preciso escrever o
receptor, a validação e os limites, e manter um formato que nenhuma outra
ferramenta lê; na prática, é o backend próprio deste trabalho sem o Collector
na frente. Pelo mesmo motivo, a operação e as regras de sucesso e de
retentativa (C6 e C7) seriam código do projeto. A maturidade (C8) também é uma
ressalva: um componente escrito e mantido pelo próprio projeto começa sem
histórico que permita julgá-lo maduro, e chegar lá pediria trabalho que vai
além do código, como formar e manter uma comunidade em torno dele. Essa linha
da tabela é avaliação deste trabalho, sem fonte externa.

O Kieker, uma
ferramenta de monitoramento com formato de registro próprio, mostra o que
acontece quando se quer reusar outra ferramenta: para alimentar o ExplorViz,
de visualização, foi preciso transformar os seus registros de fluxo de
controle em \textit{traces} do OpenTelemetry, exportados por gRPC ou Zipkin
\citep{reichelt2024kieker}. Ler essa ponte como um custo do formato próprio é
interpretação deste trabalho, não do artigo.

A segunda alternativa é o \textbf{Prometheus com o Pushgateway}. O Prometheus
coleta métricas puxando-as dos processos monitorados, por HTTP, e a visão
geral do projeto indica um intermediário, o Pushgateway, para trabalhos de
vida curta \citep{prometheus_overview}, e um comando do kw, na leitura deste
trabalho, é desse tipo. A documentação, porém, aponta dois problemas. O
Pushgateway recusa envios que tragam carimbo de tempo, de modo que o momento
do evento se perde (C3) \citep{pushgateway_readme}. Além disso, a documentação
diz que o uso válido do Pushgateway é, em geral, guardar o resultado de um
trabalho em lote de um serviço, e que ele nunca esquece as séries recebidas
\citep{prometheus_pushing}, o que sugere que muitas instalações independentes,
cada uma com seus eventos, não são o caso previsto; e, como cada atributo novo
vira um rótulo novo e cria séries novas, uma fonte nova pesa sobre esse
acúmulo (C2).

No Pushgateway, as métricas ficam agrupadas por uma chave, e um envio por PUT
substitui todas as métricas do grupo \citep{pushgateway_readme}, então o
cliente teria de reenviar o estado do grupo, não só os eventos novos (C4). O
filtro de rótulos fica no Prometheus que coleta, depois que o dado chegou
inteiro ao Pushgateway (C1) \citep{prometheus_config_v3_16_0}. O Pushgateway
não guarda os dados em disco por padrão \citep{pushgateway_readme}, e nada
nele confirma que um envio foi gravado pelo Prometheus que o coleta (C7). São,
por fim, dois serviços a operar (C6).

A terceira alternativa é uma plataforma de análise de produto, o
\textbf{PostHog}, que substituiria o caminho inteiro, do recebimento à
consulta. É código aberto, sob a licença MIT fora do diretório \texttt{ee/}
(C5) \citep{posthog_license}, e aceita eventos por uma API própria e
\textit{logs} por OTLP \citep{posthog_capture, posthog_logs}; escolher OTLP,
portanto, não fecha a porta a uma plataforma como essa. Seu ponto forte é o
que o C9 pede: o armazenamento e a consulta vêm prontos, com mais recursos que
os do backend próprio.

O PostHog também filtra propriedades dos eventos
imediatamente antes de gravá-los \citep{posthog_realtime_transformations}, mas
no servidor, depois do envio. Nas páginas consultadas, nenhuma diz se esse
recurso funciona na instalação própria, e a lista do que se perde ao hospedar
o PostHog por conta própria inclui os \textit{data pipelines}, seção onde ele
está documentado \citep{posthog_selfhost_disclaimer}. A API de captura exige
um identificador da pessoa em cada evento, enquanto o kw usaria só um
identificador aleatório de instalação (C1), e responde 200 mesmo a eventos que
não grava \citep{posthog_capture}; só a versão mais nova da API diz, evento a
evento, o que foi gravado (C7) \citep{posthog_capture_v1}. Contra ele pesa o custo
de operação (C6): a instalação própria pede o equivalente a uma máquina com
4 vCPU e 16~GB de memória \citep{posthog_selfhost}, para uma escala de cerca
de um envio por instalação por dia (Capítulo~\ref{cap:solucao}), e não tem
suporte oficial (C8) \citep{posthog_selfhost}.

\begin{table}[!ht]
  \centering
  \newcommand*{\cA}[1][]{A#1}
  \newcommand*{\cR}[1][]{\cellcolor{black!10}R#1}
  \newcommand*{\cN}[1][]{\cellcolor{black!25}N#1}
  \newcommand*{\cQ}[1][]{\textbf{?}#1}
  \newcommand*{\cNA}[1][]{n/a#1}
  \newcommand*{\nosso}{\textsuperscript{*}}
  \small
  \setlength{\tabcolsep}{3pt}
  \begin{threeparttable}
    \caption{Alternativas consideradas, avaliadas pelos critérios da
      Tabela~\ref{tab:criterios}.}
    \label{tab:alternativas}
    \begin{tabular}{@{}>{\raggedright\arraybackslash}p{3.9cm}
                      *{9}{>{\centering\arraybackslash}p{0.92cm}}@{}}
      \toprule
       & \makecell{C1\\[-2pt]\scriptsize filtro} & \makecell{C2\\[-2pt]\scriptsize extens.}
       & \makecell{C3\\[-2pt]\scriptsize evento} & \makecell{C4\\[-2pt]\scriptsize cliente}
       & \makecell{C5\\[-2pt]\scriptsize aberto} & \makecell{C6\\[-2pt]\scriptsize custo}
       & \makecell{C7\\[-2pt]\scriptsize entrega} & \makecell{C8\\[-2pt]\scriptsize matur.}
       & \makecell{C9\\[-2pt]\scriptsize reuso} \\
      \midrule
      OpenTelemetry (OTLP e Collector), a escolha
        & \cA[\nosso] & \cA[\nosso] & \cA & \cA[\nosso] & \cA
        & \cR & \cR & \cR & \cA \\
      Endpoint HTTP/JSON próprio
        & \cR[\nosso] & \cA[\nosso] & \cA[\nosso] & \cA[\nosso] & \cA[\nosso]
        & \cR[\nosso] & \cR[\nosso] & \cR[\nosso] & \cN[\nosso] \\
      Prometheus com Pushgateway
        & \cR & \cR & \cN & \cR & \cA
        & \cR & \cN & \cA & \cA \\
      PostHog (plataforma de análise)
        & \cR & \cA & \cA & \cA & \cR
        & \cN & \cR & \cR & \cA \\
      \bottomrule
    \end{tabular}
    \begin{tablenotes}
      \footnotesize
      \item A: atende; R: atende com ressalva; N: não atende; *: avaliação
        deste trabalho, sobre o design ou o repositório, sem fonte externa
        (Capítulo~\ref{cap:solucao}).
        Critérios na Tabela~\ref{tab:criterios}.
      \item Justificativas na Seção~\ref{sec:alternativas}.
    \end{tablenotes}
  \end{threeparttable}
\end{table}


\section{Precedentes em ferramentas de linha de comando}
\label{sec:precedentes}

O Capítulo~\ref{cap:motivacao} citou quatro ferramentas que já enfrentaram a
questão do kw: a cadeia de ferramentas do Go, o Debian Popularity Contest
(popcon), o Homebrew e a interface de linha de comando do .NET. Todas rodam
como processos curtos na máquina de quem as usa, sem um serviço próprio sempre
ligado, e precisam levar dados de muitas máquinas para um lugar central. Aqui
interessa como cada uma pede consentimento, envia, guarda o estado do envio e
identifica a máquina, não o que coleta, que não é critério para o kw. A
Tabela~\ref{tab:precedentes}, ao fim da seção, resume as quatro e, na última
linha, como o kw trata cada um desses pontos.

A cadeia de ferramentas do \textbf{Go} guarda contadores na máquina,
agrega-os em relatórios semanais e, por padrão, nunca os envia
\citep{go_telemetry_doc}. O design
publicado por Russ Cox em fevereiro de 2023 era opt-out
\citep{cox_telemetry_design_2023}; dezesseis dias depois, após debate público,
passou a opt-in \citep{coxsw2023optin}, e é assim que a telemetria saiu no
Go 1.23 \citep{go_release_1_23}. Os relatórios enviados, diz o design, não
incluem identificador de nenhum tipo
\citep{cox_telemetry_design_2023}. O relatório vai por POST, em JSON; só
depois de uma resposta 200 sai da fila local, e uma resposta 4xx o descarta
\citep{golang_telemetry}. Na discussão pública do design, Cox disse entender
o OpenTelemetry como organizado em torno de um fluxo de eventos e escreveu
que evitar esse fluxo é uma das formas de reduzir o impacto sobre a
privacidade. Um desenvolvedor do OpenTelemetry concordou, disse que o SDK de
Go poderia ser adaptado, mas que não valia o esforço, e que o OpenTelemetry
poderia ser um bom lugar para uma telemetria assim
\citep{go_discussion_58409}.

O \textbf{popcon} coleta dados de uso para orientar o próprio projeto, e o
artigo do kw já o toma como referência (Seção~\ref{sec:dados-de-uso}). A
participação é opt-in, com \enquote{não} como resposta padrão na instalação.
Cada máquina recebe um identificador aleatório de 128 bits, e um novo envio
sobrescreve o último de mesmo identificador; o servidor publica todo dia um
resumo que junta os envios e não traz os identificadores
\citep{popcon_package_1_78}. O próprio popcon avisa que nenhum envio é
completamente anônimo \citep{popcon_package_1_78}, ponto que o
Capítulo~\ref{cap:privacidade} retoma.

O \textbf{Homebrew} mostra um aviso antes de ativar a coleta, que pode ser
desligada, e a carga não tem identificador de usuário nem campo de endereço IP
\citep{brew_analytics}. O envio é um \texttt{curl} por HTTPS em processo
destacado, com tempo máximo de 3 segundos \citep{brew_analytics_code}, para
não atrasar o comando. As falhas são silenciosas \citep{brew_analytics}, e não
há fila nem estado de envio em disco \citep{brew_analytics_code}.

A \textbf{interface de linha de comando do .NET} é a única das quatro que adota o
OpenTelemetry, e ainda só no ramo principal e em versões de prévia. A troca
entrou em abril de 2026 \citep{dotnet_sdk_pr_53181} e está nas prévias 11.x a
partir da quarta; desde a quinta, o envio por OTLP depende de uma variável de
ambiente \citep{dotnet_sdk_telemetryclient_rc3}. As versões estáveis 10.x
ainda usam o cliente anterior, que guarda os envios numa fila em disco e os
reenvia quando o erro permite \citep{dotnet_sdk_v10_0_404}. A telemetria é
opt-out, ligada por padrão nas distribuições da Microsoft e desligável por
variável de ambiente; o identificador da máquina é o SHA-256 do endereço MAC,
e os agregados são publicados sob licença Creative Commons
\citep{dotnet_cli_telemetry}. É, portanto, um projeto grande da mesma
categoria levando a sua telemetria para o OpenTelemetry, ainda fora das
versões estáveis, e usando o SDK, que na plataforma .NET é nativo.

O \textbf{kw} usa o opt-in, como o Go e o popcon, um identificador aleatório
de instalação, como o popcon, e o envio em segundo plano, sem bloquear o
comando, como o Homebrew. Ele se afasta do Go na forma como os dados são
agregados. No Go, a agregação acontece no cliente, e só o relatório semanal,
já resumido, sai da máquina. No kw, cada evento sai como um registro, e as
agregações são feitas no backend, por consulta, o que deixa qualquer atributo
novo disponível para agrupar sem mudar o cliente (Capítulo~\ref{cap:solucao});
o custo dessa escolha está na Seção~\ref{sec:limites-otel}. Ao contrário do
Homebrew, o kw só marca um registro como enviado depois da confirmação do
backend. Ao contrário do .NET, não envia nada antes do opt-in, não deriva o
identificador de um dado da máquina e, escrito em Bash, não usa SDK.

\begin{table}[!ht]
  \centering
  \footnotesize
  \setlength{\tabcolsep}{4pt}
  \begin{threeparttable}
    \caption{Como quatro ferramentas de linha de comando pedem o
      consentimento, enviam os dados, tratam o envio e identificam a
      máquina, e como o kw trata os mesmos pontos.}
    \label{tab:precedentes}
    \begin{tabular}{@{}>{\raggedright\arraybackslash}p{1.7cm}
                      >{\raggedright\arraybackslash}p{3.0cm}
                      >{\raggedright\arraybackslash}p{3.0cm}
                      >{\raggedright\arraybackslash}p{3.4cm}
                      >{\raggedright\arraybackslash}p{2.9cm}@{}}
      \toprule
      Ferramenta & Consentimento & Como e para onde envia & Estado de envio
        & Identificador \\
      \midrule
      Go
        & Opt-in: o modo padrão, \texttt{local}, guarda os dados e nunca envia
        & Relatório semanal em JSON, por POST, para
          \texttt{telemetry.go.dev}
        & Após a resposta 200, o relatório vai para a pasta
          \texttt{upload/}; um erro 4xx o apaga
        & Nenhum: os relatórios não trazem identificador de nenhum tipo,
          segundo o design de 2023 \\
      \addlinespace
      Debian popcon
        & Opt-in: a resposta padrão na instalação é ``não''
        & POST por HTTP para \texttt{popcon.\allowbreak debian.org}
        & O servidor substitui o envio anterior da mesma máquina pelo novo
        & UUID aleatório de 128 bits por máquina \\
      \addlinespace
      Homebrew
        & Mostra um aviso antes de ativar e pode ser desligada
        & \texttt{curl} em processo destacado, por HTTPS, com tempo máximo de
          3~s
        & Nenhum estado de envio nem fila própria; as falhas são silenciosas
        & Nenhum: o envio não traz identificador de usuário nem campo de
          endereço IP \\
      \addlinespace
      CLI do .NET
        & Ligada por padrão nas distribuições da Microsoft; \mbox{opt-out} por
          variável de ambiente
        & Azure Monitor; nas prévias da versão 11, também OTLP, que depende de
          variável de ambiente a partir da \textit{preview}~5
        & Na versão 10, anterior ao OpenTelemetry: fila em disco e reenvio
          quando o erro permite
        & \textit{Hash} SHA-256 do endereço MAC \\
      \midrule
      kw
        & Opt-in; enquanto não há resposta, o estado é \texttt{ask} e nada é
          enviado; o usuário pode excluir campos do envio
          (\texttt{exclude\_fields}), vazio por padrão
        & Um registro por evento, com o tempo em granularidade de semana, OTLP/HTTP
          com JSON, em segundo plano, sem bloquear o comando, para um
          Collector e daí ao backend
        & Cada evento é marcado como enviado só depois da confirmação do backend
        & UUID aleatório por instalação, criado no opt-in \\
      \bottomrule
    \end{tabular}
    \begin{tablenotes}
      \footnotesize
      \item Justificativas na Seção~\ref{sec:precedentes}. Linha do kw:
        avaliação deste trabalho sobre o repositório
        (Capítulo~\ref{cap:solucao}).
    \end{tablenotes}
  \end{threeparttable}
\end{table}


\section{Limites da escolha}
\label{sec:limites-otel}

Reusar o OpenTelemetry tem um preço, e o primeiro está no Collector, que está
na versão 0.x; e a política de versões do projeto avisa que, nessa faixa, nada
é garantido estável \citep{otel_collector_v0_121_0}. O receptor de OTLP e o
exportador por HTTP, por onde o dado entra e sai, são estáveis para
\textit{logs}; os três \textit{processors} do caminho, o de limite de memória,
o de recurso e o de atributos, estão em beta. As mudanças incompatíveis das
versões 0.119.0 a 0.121.0, esta última a usada, não tocaram nenhum desses
componentes \citep{otel_collector_v0_121_0}. O Collector é também mais um
serviço a operar, ao lado do backend (C6). A versão fica fixada, e cada
atualização precisa ser testada (C8).

O segundo preço está no sinal e no formato. O modelo de dados de \textit{logs}
e o protocolo são estáveis \citep{otel_logs_data_model, otlp_spec_1_11_1}. O
modelo foi declarado estável em março de 2022, sem a API e os SDKs de
\textit{logs} \citep{otel_spec_pr_2387}, e o suporte a \textit{logs} no SDK de
Python ainda está em desenvolvimento \citep{otel_python_docs}; como o caminho
não usa SDK, importam o modelo e o protocolo. O JSON do OTLP é a tradução do
protobuf e herda as regras dele: ao emitir, inteiros de 64 bits, como os
carimbos de tempo em nanossegundos, saem como texto decimal, enumerações como
inteiros e identificadores em hexadecimal \citep{otlp_spec_1_11_1}. O cliente
tem de conhecer essas regras, e o corpo fica maior que o do protobuf
(Capítulo~\ref{cap:solucao}). O protocolo também não garante entrega única: a
especificação diz que reenviar para garantir a entrega pode gerar duplicatas,
e trata isso como uma troca aceita de propósito \citep{otlp_spec_1_11_1}. A
confirmação depois da gravação e o descarte de repetidos ficam com o backend
próprio (C7).

O terceiro preço é o que fica guardado no backend. O Go envia contadores agregados
por semana, não um fluxo de eventos (Seção~\ref{sec:precedentes}). Russ Cox,
autor do design, escreveu que evitar um fluxo de eventos é uma das formas
importantes de a telemetria do Go reduzir o impacto sobre a privacidade, e
disse entender o OpenTelemetry como construído em torno de um fluxo de
eventos \citep{go_discussion_58409}.

Este trabalho envia um registro por evento e
agrega no backend, pelos motivos de extensibilidade do
Capítulo~\ref{cap:solucao}. A redução de impacto vem de passos feitos no
cliente, antes do envio: o consentimento, a exclusão dos campos que o usuário
escolher, a granularidade do tempo reduzida à semana e o identificador
aleatório de cada instalação; o backend reduz de novo a granularidade do
tempo no que recebe, como segunda barreira. O custo é que o backend guarda eventos individuais, e quem o
opera vê mais do que um contador semanal mostraria.

O peso de guardar eventos individuais, porém, depende de como a ferramenta é
usada. A cadeia de ferramentas do Go é a base de quem programa na linguagem; o
kw é chamado de tempos em tempos durante o trabalho no kernel, para compilar,
instalar ou preparar o envio de \textit{patches}. Um registro a cada chamada
do kw é, na leitura deste trabalho, um fluxo de eventos modesto. Como o tempo
de cada registro sai com granularidade de semana, o fluxo também revela pouco sobre
quando o usuário trabalha, embora o trabalho no kernel deixe rastros públicos e
datados, nos \textit{commits} e nas listas de e-mail, aos quais esses
registros ainda podem ser ligados pela semana. Dois pontos ficam em aberto: a
estabilidade dos \textit{processors} em beta, que precisa ser acompanhada, e o
anonimato dos eventos individuais quando o backend reunir os de muitas
instalações, que o Capítulo~\ref{cap:privacidade} retoma.